\cleardoublepage
\chapter*{Afterword}
\addcontentsline{toc}{chapter}{Afterword}
\thispagestyle{empty}

A mathematical book usually ends at a point where mathematics itself does not.

The final theorem may have been proved, the last example may have been examined, and the notation may have reached a temporary resting place. Yet none of these provides a genuine conclusion. Every definition introduced in these pages admits further abstraction. Every theorem belongs to a larger theory. Every proof suggests questions about which assumptions are essential, which arguments can be generalized, and which structures remain undiscovered.

For this reason, the end of these notes should not be understood as the completion of a mathematical journey. It marks only the point at which one particular route has been recorded.

This manuscript began as an attempt to organize what I had learned. Like many students encountering university mathematics, I first met its subjects as separate courses. Analysis concerned limits and continuity. Linear algebra concerned vectors and matrices. Abstract algebra concerned groups, rings, and fields. Topology concerned open sets and continuous maps. Probability concerned random variables and distributions. Each subject appeared to possess its own terminology, techniques, and standards of argument.

Only gradually did their common structure become visible.

The same questions returned in different forms. What are the objects under consideration? Which relations between them matter? What transformations preserve the relevant structure? What assumptions are actually required for a conclusion? When can a local property be extended globally? When does an algebraic description correspond to a geometric one? When does convergence preserve the operations we wish to perform?

Once these questions are recognized, mathematics no longer appears as a collection of independent territories. It becomes a network of recurring ideas expressed through different languages.

Completeness appears first as a property of number systems and metric spaces, but later governs the existence of limits, fixed points, and solutions. Compactness begins as a topological condition, yet repeatedly converts local information into global control. Symmetry is encoded algebraically through groups, geometrically through transformations, and analytically through invariance. Linearity, initially encountered in systems of equations, becomes a principle for decomposing complicated objects into manageable components. Measure extends elementary notions of length and area, while probability gives those same structures a new interpretation.

To study mathematics is therefore not merely to collect results. It is to learn to recognize these patterns of thought as they reappear.

Writing these notes made that lesson more demanding. It is possible to believe that one understands a theorem while reading it. It is much harder to explain why its hypotheses are natural, reconstruct its proof without hidden steps, and place it within the surrounding theory. Exposition exposes the difference between familiarity and understanding. A sentence that seems harmless may conceal an undefined term. An ``obvious'' implication may require a lemma. A standard proof may depend on a convention established many pages earlier.

In this sense, writing became part of learning rather than a record produced after learning was complete.

The imperfections of this manuscript belong to the same process. Its subjects are unequal in scale, and no single volume can treat them with equal depth. Some chapters lead close to the frontier of a theory; others establish only the language required to begin. Certain arguments have been included because they illuminate a general method, while other important results have been omitted because their proper treatment would require a separate book. The selection reflects not a final map of mathematics, but the path by which one student attempted to see its landscape more clearly.

That path will inevitably change.

Definitions that once appeared artificial may later seem unavoidable. Proofs that once seemed difficult may become examples of a familiar technique. Conversely, statements that appear settled may reveal deeper questions when viewed from another subject. Mathematical maturity does not eliminate difficulty; it changes the level at which difficulty is encountered.

A beginning student asks how to prove a theorem. A more experienced student also asks why this theorem was formulated, how its assumptions were discovered, whether its conclusion is optimal, and what larger structure makes the proof possible. Progress is measured not only by the number of problems one can solve, but by the quality of the questions one learns to ask.

There is also a particular discipline in accepting incompleteness.

One may be tempted to postpone writing until every detail is mastered, or to delay sharing a set of notes until every chapter has reached the same degree of polish. But complete mastery is not a state that mathematics readily permits. Each answer enlarges the boundary of what remains unknown. Waiting for finality would mean never beginning.

The more honest goal is not perfection, but continued correction: to state assumptions more clearly, replace vague intuition with precise argument, discover counterexamples to careless claims, and revise explanations when a better structure becomes visible. Rigor is not achieved once and preserved automatically. It is a habit of returning to what has been written and asking whether it truly says what was intended.

The reader should approach these pages in the same spirit.

Do not treat the proofs as monuments that can only be observed from a distance. Take them apart. Reconstruct them. Search for shorter arguments and stronger conclusions. Ask what happens when an assumption is removed. Draw examples that fit a definition and counterexamples that nearly fit it. Compare different formulations of the same idea. A proof becomes part of one's mathematical understanding only after it can be recreated, adapted, and questioned.

Nor should these notes be read as an endpoint. Their natural continuation lies in the books, lectures, problems, and conversations from which mathematics is actually learned. Real analysis leads toward functional analysis, partial differential equations, and dynamical systems. Algebra leads toward representation theory, commutative algebra, and algebraic geometry. Topology leads toward manifolds, homotopy, and geometry. Measure and probability lead toward stochastic processes, statistics, and mathematical physics. The boundaries between these directions are themselves increasingly artificial.

The structure described in the preface as a mathematical edifice is therefore not a finished skyscraper. It is closer to a city under permanent construction. Some foundations are ancient; some towers are new. Different districts develop according to different needs, yet bridges continually appear between them. Old structures are strengthened, extended, or rebuilt when new perspectives demand it.

To enter mathematics is not simply to admire this city from above. It is to learn how to work within it.

One begins with small tasks: understanding a definition, checking an example, completing a proof. Over time, these acts develop into a way of thinking. Precision becomes less a formal obligation than a form of intellectual honesty. Abstraction becomes less an escape from concrete reality than a method for discovering what many concrete situations share. Proof becomes not merely a certificate of truth, but an explanation of why a statement could not have been otherwise under the given assumptions.

This is perhaps the deepest lesson I have taken from the process of writing these notes. Mathematics is not defined solely by the certainty of its conclusions. It is defined by the care with which it moves from one idea to the next.

The manuscript ends here, but its purpose extends beyond its final page. It will have succeeded if it encourages the reader to continue independently: to question established arguments, to seek connections between distant subjects, and eventually to construct explanations that are clearer than those found here.

A preface looks forward to a path that has not yet been travelled. An afterword looks back and discovers that the path has changed the traveller.

What lies ahead is larger than what has been recorded.

\begin{flushright}
Jinshuo Li \\
Shanghai Jiao Tong University \\
\end{flushright}

\cleardoublepage